% !TeX encoding = UTF-8
\documentclass[variant=apuntes,cover=institutional,mode=screen]{ua-engineering-v6-article}
% !TeX encoding = UTF-8
\UASetUniversity{Universidad Austral}
\UASetFaculty{Facultad de Ingeniería}
\UASetDepartment{Departamento de Computación}
\UASetCampus{Pilar}
\UASetCareer{Ingeniería Informática}
\UASetCommission{Comisión A}
\UASetProfessor{Equipo docente}
\UASetSemester{Segundo cuatrimestre 2026}
\UASetCourse{Sistemas Distribuidos}
\UASetDate{2026}


\UASetClassNumber{14}
\UASetClassTitle{Incidentes y cultura de confiabilidad}
\UASetDocKind{Apuntes}

\title{Clase 14 --- Incidentes, postmortems, aprendizaje organizacional y cultura de confiabilidad}
\author{\UAProfessor}
\date{\UADate}

\begin{document}
\maketitle

\section{Introducción}

En sistemas distribuidos complejos, los incidentes no son anomalías imposibles ni evidencia automática de incompetencia. Son eventos esperables en un entorno donde interactúan:

\begin{itemize}
\item fallas parciales;
\item latencias variables;
\item decisiones de diseño imperfectas;
\item carga imprevisible;
\item errores humanos;
\item acoplamientos organizacionales y técnicos.
\end{itemize}

La cuestión clave no es solo prevenir incidentes, sino desarrollar la capacidad de:
\begin{itemize}
\item detectarlos;
\item responder con claridad;
\item mitigar daño;
\item aprender de forma acumulativa.
\end{itemize}

\begin{quote}
La confiabilidad madura no consiste en ausencia total de incidentes, sino en aprendizaje operativo sistemático a partir de ellos.
\end{quote}

\section{Qué es un incidente}

Un incidente es un evento que degrada, o amenaza con degradar, una propiedad relevante del servicio. Esto puede incluir:

\begin{itemize}
\item caída o indisponibilidad;
\item aumento severo de latencia;
\item corrupción o inconsistencia visible;
\item pérdida de throughput;
\item impacto sobre SLOs;
\item exposición de un riesgo de seguridad o integridad.
\end{itemize}

La gravedad de un incidente no se define solo por su causa técnica, sino por su impacto y por el riesgo que implica para el servicio y los usuarios.

\section{Gestión de incidentes}

\subsection{Fases típicas}

Una gestión madura de incidentes suele incluir:

\begin{enumerate}
\item detección;
\item clasificación y severidad;
\item respuesta inicial;
\item mitigación;
\item comunicación;
\item recuperación;
\item análisis posterior.
\end{enumerate}

\subsection{Detección}

La detección puede provenir de:
\begin{itemize}
\item alertas automáticas;
\item dashboards;
\item anomalías observadas por operadores;
\item reportes de usuarios;
\item fallas en procesos downstream.
\end{itemize}

Un punto importante es que muchos incidentes son detectados tarde no porque sean técnicamente complejos, sino porque el sistema carece de señales adecuadas o porque las señales existen pero no son accionables.

\subsection{Clasificación}

Clasificar un incidente implica estimar:
\begin{itemize}
\item impacto;
\item alcance;
\item urgencia;
\item riesgo de escalamiento;
\item relación con SLOs y error budget.
\end{itemize}

Esto permite priorizar correctamente y evitar dos errores comunes:
\begin{itemize}
\item tratar todo como emergencia máxima;
\item subestimar incidentes que se amplifican con rapidez.
\end{itemize}

\section{Roles durante el incidente}

En incidentes de cierta complejidad, es útil diferenciar roles:

\begin{itemize}
\item \textbf{Incident Commander}: coordina, prioriza y mantiene foco;
\item \textbf{Investigadores}: siguen hipótesis técnicas;
\item \textbf{Comunicador}: mantiene informados a stakeholders;
\item \textbf{Scribe / Timeline Keeper}: registra decisiones y eventos.
\end{itemize}

\subsection{Razón}

Esto reduce confusión, evita que todos hagan lo mismo y disminuye la probabilidad de que la coordinación falle justo cuando más se necesita.

\section{Respuesta vs aprendizaje}

Durante el incidente, la prioridad es:
\begin{itemize}
\item contener daño;
\item restaurar servicio;
\item preservar evidencia útil.
\end{itemize}

No es el momento de discutir exhaustivamente culpa, arquitectura global o responsabilidades de largo plazo. Esas preguntas pertenecen al análisis posterior.

\section{Postmortems}

\subsection{Definición}

Un postmortem es el análisis estructurado posterior a un incidente. Su propósito no es “cerrar un ticket”, sino producir entendimiento organizacional acumulable.

\subsection{Contenido mínimo}

Un buen postmortem debería incluir:

\begin{itemize}
\item resumen ejecutivo;
\item impacto observado;
\item detección;
\item timeline;
\item respuesta y mitigación;
\item análisis causal;
\item factores contribuyentes;
\item acciones de mejora;
\item dueños y fechas de seguimiento.
\end{itemize}

\subsection{Timeline}

El timeline es especialmente importante porque evita reconstrucciones retrospectivas vagas y obliga a describir:
\begin{itemize}
\item qué pasó;
\item cuándo se detectó;
\item qué hipótesis se siguieron;
\item qué decisiones se tomaron;
\item qué mitigación cambió efectivamente el curso del incidente.
\end{itemize}

\section{Blameless Postmortem}

\subsection{Qué significa}

Un postmortem \emph{blameless} no significa que “nadie sea responsable” o que “todo da igual”. Significa que el análisis evita convertir el castigo individual en sustituto del entendimiento sistémico.

\subsection{Por qué importa}

Si la organización responde a incidentes buscando culpables como objetivo principal, aparecen efectos dañinos:

\begin{itemize}
\item ocultamiento de errores;
\item menor reporte temprano;
\item explicaciones superficiales;
\item aprendizaje empobrecido;
\item miedo a experimentar o escalar problemas.
\end{itemize}

\subsection{Lectura correcta}

Blameless no elimina responsabilidad profesional. Lo que elimina es la ilusión de que la explicación del incidente se agota en un individuo que “cometió un error”.

\section{Causa raíz y factores contribuyentes}

\subsection{Problema de la causa única}

En sistemas complejos, reducir un incidente a “la causa raíz” puede ser útil como simplificación, pero a menudo es insuficiente o engañoso. Más realista es distinguir:

\begin{itemize}
\item disparador inmediato;
\item condiciones habilitantes;
\item debilidades de diseño;
\item fallas de observabilidad;
\item fallas de mitigación;
\item fallas de coordinación.
\end{itemize}

\subsection{Ejemplo}

Un incidente visible como “API 500” puede involucrar:
\begin{itemize}
\item latencia alta en DB;
\item retries en cascada;
\item breaker ausente;
\item pools compartidos;
\item alertas tardías;
\item runbook insuficiente.
\end{itemize}

La narrativa correcta no es “la base estuvo lenta”, sino la cadena completa que transformó esa lentitud en una degradación sistémica.

\section{Acciones de mejora}

\subsection{Buenas acciones}

Las acciones posteriores al postmortem deben ser:
\begin{itemize}
\item concretas;
\item verificables;
\item con owner;
\item con horizonte temporal;
\item orientadas a reducir repetición o impacto.
\end{itemize}

\subsection{Tipos}

\begin{itemize}
\item cambios de código;
\item mejoras de observabilidad;
\item ajustes de umbrales;
\item nuevos tests de falla;
\item revisiones arquitectónicas;
\item cambios en runbooks;
\item entrenamiento del equipo.
\end{itemize}

\subsection{Malas acciones}

\begin{itemize}
\item “tener más cuidado”;
\item “mejorar comunicación” sin mecanismo concreto;
\item “vigilar más”;
\item acciones sin responsable.
\end{itemize}

\section{Aprendizaje organizacional}

\subsection{Qué significa aprender}

Aprender de incidentes no es solo corregir el bug puntual. Es cambiar la capacidad futura de la organización para:
\begin{itemize}
\item detectar antes;
\item responder mejor;
\item diseñar con más criterio;
\item reducir superficie de repetición.
\end{itemize}

\subsection{Señal de madurez}

Una organización madura no se pregunta solo:
\begin{quote}
“¿Cómo arreglamos este incidente?”
\end{quote}

Se pregunta también:
\begin{quote}
“¿Qué hizo posible que esto ocurriera, se expandiera y tardara tanto en entenderse?”
\end{quote}

\section{Cultura de confiabilidad}

La cultura de confiabilidad no es únicamente una colección de herramientas. Incluye:

\begin{itemize}
\item visibilidad de incidentes;
\item disciplina de postmortems;
\item ownership claro;
\item aprendizaje sin ocultamiento;
\item equilibrio entre cambio y estabilidad;
\item valoración explícita de la operación como trabajo técnico serio.
\end{itemize}

\subsection{Anti-patrones culturales}

\begin{itemize}
\item tratar incidentes como vergüenza a ocultar;
\item celebrar solo velocidad de entrega;
\item ignorar deuda operacional;
\item escribir postmortems vacíos;
\item no cerrar el loop de acciones.
\end{itemize}

\section{Relación con SLOs y error budgets}

Los incidentes se interpretan mejor cuando existen:
\begin{itemize}
\item SLI claros;
\item SLOs explícitos;
\item error budgets definidos.
\end{itemize}

Eso permite responder preguntas como:
\begin{itemize}
\item ¿qué tan grave fue el impacto?;
\item ¿consumió cuánto budget?;
\item ¿qué cambios de ritmo exige?;
\item ¿qué trade-off entre cambio y estabilidad debemos revisar?
\end{itemize}

\section{Modelo del curso}

Podemos expresar el problema como:

\[
D = (P, F, G, S, T)
\]

Por ejemplo:
\begin{itemize}
\item $P$: responder y aprender de incidentes en un sistema distribuido;
\item $F$: fallas parciales, ruido observacional, coordinación imperfecta, decisiones bajo presión;
\item $G$: recuperación razonable y aprendizaje organizacional efectivo;
\item $S$: incident management + postmortems + ownership + cultura blameless;
\item $T$: tiempo, disciplina, exposición de debilidades y costo organizacional de mejorar.
\end{itemize}

\section{Conclusión}

Los incidentes no son solo eventos técnicos. Son también eventos organizacionales que revelan:
\begin{itemize}
\item límites de la arquitectura;
\item límites de la observabilidad;
\item límites del proceso operativo;
\item límites de la cultura.
\end{itemize}

\begin{uakeybox}
La confiabilidad madura no consiste en “tener menos incidentes” por arte de magia, sino en convertir cada incidente en una oportunidad real de aumentar la capacidad del sistema y de la organización.
\end{uakeybox}


\section{Síntesis razonada de la clase}
Esta unidad debe leerse como parte de una cadena de decisiones de diseño y no como un catálogo de mecanismos. Para estudiar el tema con profundidad conviene reconstruir cuatro preguntas: \emph{qué problema aparece}, \emph{qué garantía se desea}, \emph{qué mecanismo la aproxima} y \emph{qué costo introduce}. En cada ejemplo debe identificarse además el modelo de falla implícito.

\begin{uakeybox}
Una respuesta técnicamente madura no termina al nombrar un algoritmo o patrón: declara sus supuestos, la propiedad que preserva y el escenario en el que deja de ser suficiente.
\end{uakeybox}

\section{Bibliografía guiada y ruta de profundización}
Para profundizar esta clase se recomienda: Google SRE Book y SRE Workbook; literatura de resilience engineering.

La lectura debe hacerse con preguntas concretas: (1) ¿qué modelo de sistema asume el autor?, (2) ¿qué propiedad demuestra o persigue?, (3) ¿qué falla o anomalía motiva el mecanismo?, y (4) ¿qué trade-off queda abierto?

\section{Conexión con el Trabajo Final Integrador}
El grupo debe señalar qué parte de su sistema utiliza los conceptos de esta unidad, justificar por qué son necesarios y documentar una alternativa descartada. Cuando el contenido todavía no corresponda a la etapa vigente, debe registrarse como hipótesis o decisión pendiente.

\section{Autoevaluación conceptual}
\begin{enumerate}
\item ¿Cuál es el problema distribuido concreto que motiva esta clase?
\item ¿Qué propiedad de seguridad y qué propiedad de progreso están en juego?
\item ¿Qué supuesto de red, tiempo o falla resulta decisivo?
\item ¿Qué trade-off aceptaría en un sistema real y cuál no aceptaría en un dominio crítico?
\item ¿Cómo observaría experimentalmente que la solución se comporta como espera?
\end{enumerate}

\section{Profundización autocontenida v9}
\subsection{Problema, límite e invariante}
El núcleo de esta clase es incident command, timeline, postmortems blameless, factores contribuyentes y aprendizaje organizacional. Para estudiarlo de manera autónoma conviene separar tres planos. Primero, el \emph{problema observable}: qué comportamiento incorrecto puede ver un cliente o un operador. Segundo, el \emph{límite del modelo}: qué información, sincronía o coordinación no está disponible. Tercero, el \emph{invariante}: qué propiedad no debe violarse aunque el sistema pierda progreso temporalmente.

El modelo de fallas relevante incluye detección tardía, coordinación confusa, causa simplificada y acciones sin owner. Una solución debe describirse como protocolo y no solo como nombre de tecnología: actores, estados, mensajes, condición de éxito, condición de aborto y estado visible después de una interrupción. La evidencia esperada es timeline, impacto, factores contribuyentes y acciones verificables.

\subsection{Método de análisis}
Ante un caso nuevo, siga esta secuencia:
\begin{enumerate}
\item Identifique actores, dato o recurso compartido y frontera de confianza.
\item Escriba una ejecución nominal y una ejecución adversarial.
\item Distinga seguridad (lo malo no ocurre) de progreso (algo bueno finalmente ocurre).
\item Declare qué supuesto permite avanzar: mayoría, deadline, sincronía parcial, operación idempotente, almacenamiento durable u otro.
\item Compare una alternativa más simple y una más fuerte; registre qué métrica o experimento haría cambiar la decisión.
\end{enumerate}

\subsection{Caso transferible}
Considere una plataforma de reservas que recibe operaciones duplicadas, pierde conectividad entre regiones y debe sostener un camino crítico sin ocultar el estado real al usuario. Aplique el mecanismo de esta clase a una única decisión concreta. Una respuesta madura no intenta resolver todo: delimita la operación, formula la garantía y explica la degradación esperada cuando el supuesto central deja de cumplirse.

\subsection{Errores de razonamiento frecuentes}
\begin{itemize}
\item confundir una propiedad con el producto que suele implementarla;
\item describir el camino feliz sin recorrer una falla intermedia;
\item afirmar ``exactly-once'', ``alta disponibilidad'' o ``consistencia'' sin definir alcance observable;
\item medir solo throughput aunque la hipótesis trate sobre semántica o recuperación;
\item presentar la complejidad añadida como beneficio sin compararla con la alternativa más simple.
\end{itemize}

\section{Lectura guiada}
Google SRE Workbook, incident response y postmortems; literatura de safety science.

Durante la lectura, registre: modelo de sistema, propiedad perseguida, contraejemplo que motiva el mecanismo, costo principal y una condición bajo la cual no lo elegiría. Estas cinco anotaciones convierten la bibliografía en una herramienta de diseño y no en una lista para memorizar.

\section{Aplicación al Trabajo Final Integrador}
El equipo debe producir un ADR breve con la decisión asociada a esta clase. Debe contener: contexto, dos alternativas, supuesto decisivo, garantía elegida, trade-off, evidencia pendiente y fecha de revisión. La evidencia mínima esperada para esta unidad es: timeline, impacto, factores contribuyentes y acciones verificables.

\section{Autoevaluación avanzada}
\begin{enumerate}
\item ¿Puedo explicar la clase sin nombrar primero una tecnología?
\item ¿Puedo construir una ejecución que viole la garantía si el mecanismo se configura mal?
\item ¿Puedo distinguir seguridad, progreso y experiencia del usuario?
\item ¿Puedo proponer una medición que valide la propiedad, no solo el rendimiento?
\item ¿Puedo defender cuándo una solución más simple sería suficiente?
\end{enumerate}

\end{document}
